% =============================================================================
% Supplementary Section S1 : Sensitivity to environmental proxy definition
% Response to Reviewer 1 point 1 and Reviewer 2 major point 2
% =============================================================================

\section*{S1. Sensitivity to the choice of environmental proxy}
\label{sec:S1}

\subsection*{Motivation}

The main analysis uses the SPARC environmental classification derived from
NED galaxy counts within a 2 Mpc projected radius, which yields four
discrete categories (void, field, group, cluster). A legitimate
methodological question is whether this
definition may be sensitive to the specific choice of radius, which
motivates testing alternative environmental proxies including
nth-nearest-neighbour density, host halo mass, and three-dimensional
density estimates.

We address this concern by performing a comprehensive sensitivity test:
for each SPARC galaxy, we compute 3D galaxy counts within fixed-radius
spheres at six different scales (1, 2, 3, 5, 7, and 10 Mpc) using the
2MASS Redshift Survey (2MRS) as an independent reference catalog. We
then construct a categorical proxy (4 quartiles) at each radius and fit
the same quadratic model used in the main analysis.

\subsection*{Cross-matching with 2MRS}

After cross-matching the 175 SPARC galaxies with the 2MRS catalog
(43{,}390 galaxies with measured distances), 169 galaxies have all required
coordinates (RA, Dec, distance) and BTFR residual measurements. Their
3D Cartesian positions are computed from heliocentric distances.

For each SPARC galaxy, we count the number of 2MRS galaxies within
a sphere of radius $R$. We then construct an environmental proxy by
classifying SPARC galaxies into four quartiles of this count, assigning
density proxy values $\{0.2, 0.4, 0.6, 0.8\}$ analogous to the original
SPARC classification scheme.

\subsection*{Results}

Table~\ref{tab:S1_sensitivity} summarizes the quadratic fits at each
radius, compared with the reference SPARC catalog-based classification.

\begin{table}[htbp]
\centering
\caption{Quadratic coefficient $a$ for different environmental proxy
definitions on the SPARC sample.}
\label{tab:S1_sensitivity}
\begin{tabular}{lrrrr}
\toprule
\textbf{Method} & $a$ & $\sigma_a$ & $\sigma$ & $p$ \\
\midrule
SPARC catalog (REF, 2 Mpc NED) & $+1.329$ & $0.253$ & $5.26$ & $<10^{-6}$ \\
2MRS quartiles, $R = 1$ Mpc & --- & --- & --- & (degenerate)\\
2MRS quartiles, $R = 2$ Mpc & $+0.263$ & $0.289$ & $0.91$ & $0.36$ \\
2MRS quartiles, $R = 3$ Mpc & $+0.449$ & $0.263$ & $1.71$ & $0.09$ \\
2MRS quartiles, $R = 5$ Mpc & $-0.062$ & $0.261$ & $0.24$ & $0.81$ \\
2MRS quartiles, $R = 7$ Mpc & $-0.062$ & $0.257$ & $0.24$ & $0.81$ \\
2MRS quartiles, $R = 10$ Mpc & $-0.788$ & $0.246$ & $3.20$ & $0.002$ \\
\bottomrule
\end{tabular}
\end{table}

\begin{figure}[htbp]
\centering
\includegraphics[width=\textwidth]{figureS1_sensitivity.png}
\caption{\textbf{Sensitivity of the U-shape detection to environmental
proxy definition.} Panel A shows the reference fit using the SPARC
catalog-based environmental classification (used in the main manuscript),
yielding a robust U-shape with $a = +1.33$ at $5.26\sigma$. Panels B-G
show fits using a 4-quartile classification derived from 2MRS galaxy
counts within 1, 2, 3, 5, 7, and 10 Mpc spheres respectively; the
quartile classification is degenerate at $R = 1$ Mpc, for which no fit
is shown. Panel H
provides a synthesis of the coefficient $a$ as a function of radius
compared with the SPARC reference, with the degenerate radius excluded.}
\label{fig:S1}
\end{figure}

\subsection*{Interpretation}

The reference SPARC catalog-based classification (panel A of
Figure~\ref{fig:S1}) and the 2MRS quartile-based 3D density
classifications (panels B-G) measure \textbf{physically different}
quantities:

\begin{itemize}
    \item \textbf{Catalog-based} (panel A): membership in documented
    virialized structures (Virgo cluster, Ursa Major group, etc.) that
    reflect the long-term gravitational and astrophysical history of
    each galaxy.
    \item \textbf{2MRS counts} (panels B-G): instantaneous 3D galaxy
    density within a fixed-radius sphere, which samples the local
    density agnostically of structure context.
\end{itemize}

A galaxy in the outskirts of the Virgo cluster has high catalog-based
environmental classification (R-dominated regime in the QO+R framework)
but may have relatively low instantaneous 3D galaxy density. These two
quantities encode different aspects of ``environment'': structural
context versus Poisson-like density.

The U-shape signal in BTFR residuals correlates with \textbf{structure
membership}, not with 3D density per se. This is qualitatively consistent with the QO+R framework, which predicts environmental effects from the
dominance of antagonistic Q (gas-philic, diffuse environments) and R
(stellar-philic, virialized clusters) fields. The complementary ALFALFA analysis (which uses density-based
classification at much larger scales over a much larger sample, see main
manuscript) supports the view that the signal is not specific to a single proxy
method.

\subsection*{Conclusion}

The U-shape detected in SPARC is strongest under the catalog-based
structural classification used in the main analysis. Fixed-aperture
3D density proxies yield weaker and scale-dependent trends, indicating
that the signal is sensitive to the environmental definition and may
trace structure membership or environmental history rather than
instantaneous local density alone. We interpret this as a physical
feature consistent with the QO+R framework, which predicts that
environmental effects emerge from the dominance of antagonistic
Q (gas-philic, diffuse environments) and R (stellar-philic, virialized
clusters) fields rather than from local number density per se.
However, this physical interpretation should be regarded as suggestive
rather than conclusive at this stage. A more definitive test would
require larger samples with consistent environmental classification
across all density regimes.